\section*{Summary}

A dark-matter particle that couples to the electron density breaks pairs in a superconductor through the loss function, and the loss function that conserves charge adds the pair-field ladder of \cref{sec:7.2} to the quasiparticle bubble and screens it with the Coulomb interaction (\cref{sec:15.1}). The absorptive part of the quasiparticle bubble is the Lindhard value times the case-I coherence integral $\mathfrak{g}_{\mathrm{I}}(\omega/\Delta_{\mathrm{BCS}})$, exactly, with $\mathfrak{g}_{\mathrm{I}} = 1.500$ at $y = 2.2$ falling to 1.051 at $y = 10$. The ladder raises that absorption by up to \SI{44}{\percent} near the edge. On the kinematic edge $q = \omega/v_{\mathrm{max}}$ of aluminum the Lindhard function is within \SI{10}{\percent} of the conserving loss function for $y \ge 7.9$, and the product $\mathfrak{g}_{\mathrm{I}}L^{\mathrm{nor}}$ within \SI{9}{\percent} for $y \ge 3$ at every momentum transfer. The bulk gap of aluminum is the measured \SI{0.18}{meV}; the \SI{0.3}{meV} of projections is a film convention. For a \SI{1}{keV} particle with a light mediator the bulk gap raises the rate by 1.23, more than any choice among the loss functions changes it.

On the kinematic edge each truncation is worst (\cref{sec:15.2}). There $\omega/qv_{\mathrm{F}} = 0.392$, and the normal dielectric function has an imaginary part 0.736 times its real part. In a clean crystal the conserving loss function, the physical one, is 0.87 times $\mathfrak{g}_{\mathrm{I}}L^{\mathrm{nor}}$ at $y = 2.2$ and 1.07 at $y = 3$ and 5. Impurity scattering changes the conserving value only near the threshold: at $y \ge 3$ the ratio stays between 1.008 and 1.091 for mean free paths from \SI{100}{\micro m} to \SI{30}{nm}, whereas at $y = 2.2$ it falls from 0.87 to 0.19. The superconducting part of a rate therefore carries three uncertainties: the gap, the relaxation rate of the film, and the edge remainder, which is two-sided, $-\SI{2}{\percent}$ to $+\SI{9}{\percent}$ about $\mathfrak{g}_{\mathrm{I}}L^{\mathrm{nor}}$ for $y \ge 3$. The energies $y < 3$, where the conserving value moves by more than \SI{10}{\percent} with the mean free path, are excluded from the signal region. The frequency dependence of the gap in strong coupling is not sized.

Real sensors are films (\cref{sec:15.3}). Nonmagnetic disorder leaves the gap unchanged, by Anderson's theorem, but collisions change the response. At momentum transfers below the inverse mean free path the normal loss function becomes the Drude tail $\omega\gamma_{\mathrm{rel}}/\omega_p^2$, fixed by the residual resistivity, 24 times the clean value for a \SI{50}{nm} aluminum film at \SI{0.90}{meV}. At the edge of a dirty film the density couples through the current, with the case-II factor, and the condensate short-circuits a slow potential, so the loss function of the superconducting film starts from zero and is a quarter of the normal one at $y = 2.2$. A film thinner than $1/q$ screens like a sheet, which enhances the structure factor by $[(\varepsilon_a + \varepsilon_b)/q_\parallel d_{\mathrm{film}}]^2$ to leading order in $q_\parallel d_{\mathrm{film}}$, 540 for the \SI{50}{nm} film in vacuum at \SI{0.90}{meV} (588 for the exact slab), and makes the rate directional. Amorphous WSi, with $k_{\mathrm{F}}\mathfrak{l}_e$ of a few, lies outside every relaxation-time form and outside mean-field BCS.

A boson absorbed at $\vect{q} \to 0$ needs a third body in a metal, and the Nambu structure of the current decides the coherence factor (\cref{sec:15.4}). The uniform current is the unit matrix in Nambu space and commutes with the mean-field Hamiltonian, so a clean superconductor does not absorb it. With impurities it carries the case-II factor of Mattis and Bardeen. With a phonon the second-order amplitude collapses exactly to the velocity difference times the phonon vertex, the pair carries the case-I factor, and the result is Allen's form, which the dark-matter literature evaluates with a Debye transport function. Near the edge phase space, not coherence, removes the phonon-assisted absorption: it reaches \SI{90}{\percent} of the normal-metal value at $11.4\,\Delta_{\mathrm{BCS}} = \SI{2.05}{meV}$ and is 0.42 of it at \SI{1}{meV}, where a dark photon needs $\zeta = 1.5\times10^{-9}$ for three events per kilogram-year in pure aluminum. In a \SI{50}{nm} film impurity-assisted absorption is $1.4\times10^6$ times larger at that mass. The spin current, which absorbs the axion-like particle, has the Nambu matrix $\boldsymbol{\eta}$ of the density and, in a clean singlet superconductor without spin--orbit coupling, breaks pairs by itself. In the superconductor the spin-current weight moves above $2\Delta_{\mathrm{BCS}}$, \SI{87}{\percent} of it below $4\Delta_{\mathrm{BCS}}$. The sum rule fixes the total weight; the coherent part near the edge carries the factor $(1 + F_1^a/3)/(1 + \lambda_{\mathrm{ep}})$, 0.70 for $F_1^a = 0$, and the line shape depends on $F_1^a$ through the form of Leggett and on a triplet interaction. At \SI{0.72}{meV} three events per kilogram-year in a bulk crystal with a mean free path longer than the size of a pair need $g_{aee} = 6.1\times10^{-13}$, and $5.4\times10^{-13}$ to $7.1\times10^{-13}$ for $|F_1^a/3| \le 0.1$; a \SI{50}{nm} film needs $1.5\times10^{-12}$. The strong-coupling factor $1/(1 + \lambda_{\mathrm{tr}})^2$ of the normal metal is known; whether it carries over to the superconducting phonon-assisted form is open.

\Cref{tab:15-nuisances} collects the truncations and inputs of the chapter with their sizes. For scattering in a bulk crystal the gap and the edge remainder dominate, and both are small above $y = 3$; for a film the mean free path and the thickness enter at least as strongly as the choice of response. For absorption the purity of the sample decides between impurity-assisted and phonon-assisted absorption by six orders of magnitude at \SI{1}{meV}, and for the axion-like particle in a clean crystal the direct pair breaking by the spin current exceeds the phonon-assisted absorption up to \SI{10.9}{meV}. How the quasiparticles that dark matter creates are counted is a question of the sensor readout, which this chapter does not treat.

\begin{table}[tbp]
  \caption{Uncertainty budget of the superconducting rates. Type in the classes of \cref{sec:8.8}: normalization, tilt (in $\omega$ or in the dark-matter mass), threshold and new support. Basis of the size: ``computed'' (next order or exact result computed), ``model'' (computed in a model), ``est.''\ (parametric estimate), ``input spread'' (spread or limitation of a measured or computed input), ``unsized'' (no size known). Sizes are for aluminum unless stated; $y = \omega/\Delta_{\mathrm{BCS}}$, and ``edge'' is the kinematic edge $q = \omega/v_{\mathrm{max}}$.}
  \label{tab:15-nuisances}
  \centering\scriptsize\setlength{\tabcolsep}{3pt}
  \begin{tabular}{@{}>{\raggedright\arraybackslash}p{0.26\linewidth}>{\raggedright\arraybackslash}p{0.12\linewidth}>{\raggedright\arraybackslash}p{0.35\linewidth}>{\raggedright\arraybackslash}p{0.09\linewidth}>{\raggedright\arraybackslash}p{0.10\linewidth}@{}}
    \toprule
    Truncation or input & Type & Size & Basis & Section \\
    \midrule
    Gap of the sensor & threshold; tilt in $y$ & bulk \SI{0.18}{meV} against the film convention \SI{0.30}{meV}: rate $\times1.23$ at \SI{1}{keV} (light mediator), $\times1.13$ (heavy mediator) & input spread & \cref{sec:15.1,sec:15.2} \\
    Edge remainder (central value $\mathfrak{g}_{\mathrm{I}}L^{\mathrm{nor}}$) & normali\-zation, two-sided & $-\SI{2}{\percent}$ to $+\SI{9}{\percent}$ for $y \ge 3$ (clean at all momenta, impurities on the edge); energies $y < 3$ excluded, where the conserving value moves by more than \SI{10}{\percent} with the mean free path & computed & \cref{sec:15.2} \\
    Relaxation rate of the film & tilt near the edge & factor 4.6 at $y = 2.2$ between the clean limit and $\mathfrak{l}_e = \SI{30}{nm}$; below \SI{6}{\percent} for $y \ge 3$ & computed & \cref{sec:15.2,sec:15.3} \\
    Collisions in the normal loss function & normali\-zation & Drude tail $6/\pi q\mathfrak{l}_e$ times the clean value for $q\mathfrak{l}_e \ll 1$: $\times24$ for the \SI{50}{nm} film at \SI{0.90}{meV}; neither limit within \SI{10}{\percent} for $0.64 < q\mathfrak{l}_e < 9.8$ & computed & \cref{sec:15.3} \\
    Film thinner than $1/q$ (sheet screening) & normali\-zation; tilt in direction & enhancement $[(\varepsilon_a + \varepsilon_b)/q_\parallel d_{\mathrm{film}}]^2$ of the structure factor to leading order in $q_\parallel d_{\mathrm{film}}$, $\times540$ for $\vect{q}$ in the plane of the \SI{50}{nm} film in vacuum at $qd_{\mathrm{film}} = 0.086$ ($\times588$ for the exact slab); depends on the film model and the media on both sides & model & \cref{sec:15.3} \\
    Strong disorder (amorphous WSi) & normali\-zation & $k_{\mathrm{F}}\mathfrak{l}_e = 1.8$--5.5; outside relaxation-time forms and mean-field BCS; no correction factor & unsized & \cref{sec:15.3} \\
    Thermal quasiparticles & new support & $\ee^{-\Delta_{\mathrm{BCS}}/T} = \ee^{-209}$ at \SI{10}{mK} & computed & \cref{sec:15.2} \\
    Energy dependence of the gap (strong coupling) in scattering & normali\-zation & size unknown; aluminum is near weak coupling, $2\Delta_{\mathrm{BCS}}/T_c = 3.54$ against 3.53 & unsized & \cref{sec:15.2} \\
    Shape of $\alpha^2F_{\mathrm{tr}}$ below a few meV (phonon-assisted absorption) & normali\-zation; tilt in $m_\phi$ & $\sigma_1^{\mathrm{S}}/\sigma_1^{\mathrm{N}} = 0.42$ ($\omega'^4$) or 0.73 ($\omega'^2$) at \SI{1}{meV}; normal-metal formula within \SI{10}{\percent} from $11.4$ or $7.4\,\Delta_{\mathrm{BCS}}$ & model & \cref{sec:15.4} \\
    Strong-coupling renormalization of phonon-assisted absorption & normali\-zation & $\times0.52$, the factor $1/(1 + \lambda_{\mathrm{tr}})^2$ of the normal metal at $\omega \ll \omega_{\mathrm{D}}$; whether it applies to the superconducting form is open & est. & \cref{sec:15.4} \\
    Impurity-assisted absorption & normali\-zation & $1.4\times10^6$ times the phonon-assisted $\sigma_1$ at \SI{1}{meV} for the \SI{50}{nm} film; follows the residual resistivity & computed & \cref{sec:15.4} \\
    Electron--electron absorption (intrinsic $\sigma_1$) & normali\-zation & a third intrinsic channel of $\sigma_1$ besides the phonon-assisted one; size unknown & unsized & \cref{sec:15.4} \\
    Spin-current absorption beyond the clean BCS model & normali\-zation; tilt in $m_\phi$ & total weight fixed by the sum rule; coherent weight times $1/(1 + \lambda_{\mathrm{ep}}) = 0.70$ included; Landau parameter $|F_1^a/3| \le 0.1$: $g_{aee}$ $-12$ to $+\SI{15}{\percent}$; triplet interaction, spin--orbit coupling and the phonon-assisted part near the edge unsized; clean form only for $\mathfrak{l}_e \gg \xi_{\mathrm{coh}}$ & unsized & \cref{sec:15.4} \\
    Scalar absorption near the edge (amplitude channel) & normali\-zation & kinetic row couples to the amplitude channel of the pair field; size unknown & unsized & \cref{sec:15.4} \\
    Interband absorption & new support & above about \SI{0.5}{eV}; taken from measured optical data & input spread & \cref{sec:15.4} \\
    \bottomrule
  \end{tabular}
\end{table}
